\documentclass[10pt,a4paper]{amsart}
\usepackage[margin=19mm]{geometry}
\usepackage{amsmath,amssymb,mathtools}
\usepackage[scr=rsfs,cal=euler]{mathalfa}
\usepackage{xfrac}
\usepackage[unicode,hidelinks,bookmarks=false]{hyperref}
\newcommand{\ie}{i.e.\ }
\newcommand{\cf}{cf.\ }
\newcommand{\resp}{respectively\ }
\newcommand{\Subsection}{\subsection}
\providecommand{\nobreakdash}{\nobreak\hbox{-}}
\setlength{\emergencystretch}{2em}
\multlinegap0pt
\usepackage{bm}
\usepackage{bbm}
\usepackage{dsfont}
\usepackage{xspace}
\usepackage{cancel}
\usepackage{mathtools}
\usepackage{xcolor}
\usepackage{amsthm}
\makeatletter
\@ifundefined{thm}{\newtheorem{thm}{Thm}}{}
\@ifundefined{definition}{\newtheorem{definition}[thm]{Definition}}{}
\@ifundefined{proposition}{\newtheorem{proposition}[thm]{Proposition}}{}
\@ifundefined{remark}{\newtheorem{remark}[thm]{Remark}}{}
\makeatother
\makeatletter
\newcommand{\N}{{\mathbb N}}  % Natural numbers
\newcommand{\X}{\vtr{X}}
\expandafter\ifx\csname notation\endcsname\relax
\newenvironment{notation}
%{\subsection*{Notation}\begin{description}}
%{\end{description}}

\fi
\newcommand{\vtr}[1]{\vec{\mathbf{#1}}}
\makeatother
\title[General Recurrence Multidimensional Zeckendorf Representatio]{General Recurrence Multidimensional Zeckendorf Representations}
\author{Jiarui Cheng, Steven J. Miller, Sebastian Rodriguez-Labastida, Tianyu Shen, Alan Sun, Garrett Tresch}
\date{Source: arXiv:2510.07237v1 (CC BY 4.0)}
\begin{document}
\maketitle

\noindent\emph{Source and licence.} General Recurrence Multidimensional Zeckendorf Representations. Version arXiv:2510.07237v1 (CC BY 4.0), distributed under the Creative Commons Attribution 4.0 International licence, as recorded in the arXiv licence metadata for this exact version and shown on the abstract page https://arxiv.org/abs/2510.07237v1. Full rights notice: licenses/arxiv-2510.07237-CC-BY-4.0.txt.\par\smallskip

\noindent\textit{Editorial scope.} Counted unit: the Proposition whose source label is \texttt{prop: general cSR} in the article (source0.tex lines 556--558, source label \texttt{prop: general cSR}), delivered together with the article's own complete original proof of it (source0.tex lines 559--659, one \texttt{proof} environment of 101 lines). No mathematics is written, repaired or re-derived: statement and proof are the article's own text line for line. Required context retained uncounted: defintion:NSR (lines 344--359); Rem:sumofcoefs bor/car (lines 472--474); prop:End complete (lines 525--527). Every definition, display and label of those lines is retained unchanged. Omitted from this package and not redistributed: 1--343, 360--471, 475--524, 528--555, 660--997. Nothing inside a retained excerpt is omitted or altered: every retained body is delivered exactly as the article's lines are written, so no omission receipt of any kind accompanies this package. Citations: the retained excerpts carry no citation command at all, so no bibliography entry of the article is transcribed and no reference is redistributed. Reference adaptation: none was needed and none was made; every reference inside the delivered unit resolves inside it, so no unresolved reference or citation is produced. Printed slips kept literal: no printed word, symbol, hypothesis or number of the counted statement or of its proof was altered. Layout and class: the article's own class is \texttt{interact} with packages including array, latexsym, anyfontsize, tikz, url, breakurl, amsmath, amsthm, graphicx, color, enumerate, xcolor, spverbatim, setspace; the wrapper is the lane's pinned \texttt{amsart} editorial wrapper at 10pt on A4, which restates the theorem-like environments the retained text needs and adds the article's own macro definitions that the retained text needs (\texttt{\textbackslash N}, \texttt{\textbackslash X}, \texttt{\textbackslash vtr}), with extra wrapper packages (bm, bbm, dsfont, xspace, cancel, mathtools, xcolor). The delivered numbering is the unit's own. Assets: no retained span contains a figure environment, a graphics-inclusion command, a PDF-page inclusion, a table or a TikZ picture; no image file is copied into this package, the package declares no asset dependency and no image file is redistributed with it. Licence: version arXiv:2510.07237v1 of this article is distributed under the Creative Commons Attribution 4.0 International licence, as recorded in the arXiv licence metadata for this exact version and shown on the abstract page https://arxiv.org/abs/2510.07237v1. A full rights notice is delivered at licenses/arxiv-2510.07237-CC-BY-4.0.txt. Characters: every retained excerpt is pure ASCII, so the delivered engine, XeLaTeX with its default Latin Modern text font, reports no missing character.
\begin{definition}
\label{defintion:NSR}
    A \textbf{$\vtr{c}$-nearly satisfying representation ($\vtr{c}$-NSR) for $\vtr{v} \in \mathbb{Z}^{k-1}$} is a sequence of nonnegative integers $(a_n)_{n=1}^\infty$ such that $\vtr{v}=\sum_{n=1}^\infty a_n \X_{-n}$ and where there exists an integer $i\in \N $ for which the following hold. 
        \begin{itemize}
            \item The sequence $(a_n)_{n=1}^\infty$ is not a $\vtr{c}$-SR.
            \item The sequence $(b_n)_{n=1}^\infty$ defined by
            
            $$b_n\ =\ \begin{cases}
                a_n-1, & \text{if }n=i; \\
                a_n, & \text{otherwise}
            \end{cases}$$

            \noindent is a $\vtr{c}$-SR.
        \end{itemize}
        
\end{definition}

\begin{remark}\label{Rem:sumofcoefs bor/car}
    Note that for a $\vtr{c}$-representation of $\vtr{v}$, say $a:=(a_n)_{n=1}^\infty$, with $G(a)=R$, if $b$ and $d$ are the $\vtr{c}$-representations of $\vtr{v}$ resulting from carrying into some $a_n$ and borrowing from some $a_n$ respectively, then $G(b)=R - \sum_{i=1}^m c_i + 1$ and $G(d)=R + \sum_{i=1}^kc_i -1$.
\end{remark}

\begin{proposition}\label{prop:End complete}
    Each end complete $\vtr{c}$ -NSR can be transformed into a $\vtr{c}$-SR by a finite sequence of carrying operations. Furthermore, by Remark \textcolor{blue}{\ref{Rem:sumofcoefs bor/car}}, this process only reduces the sum of the coefficients.
\end{proposition}

\begin{proposition}\label{prop: general cSR}
Every $\vtr{c}$-NSR, where $\vtr{c}$ is a weakly decreasing vector and $c_k = 1$, can be transformed into a $\vtr{c}$-SR by a finite number of borrowing and carrying operations.
\end{proposition}

\begin{proof}

Recall that if $\vtr{c}$ is a weakly decreasing vector then $c_{i+1} \leq c_i$ for all $i \in \{1,\ \ldots,\ k-1 \}$. 
Throughout this proof,  the symbol $a^{i}$ represents the $i$\textsuperscript{th} iteration of a process that consists of a finite number of borrows and carries.
Suppose that $a^{0}:=( a_n^{0} )_{n=1}^{\infty}$ is a $\vtr{c}$-NSR and let $c:=\sum_{i=1}^k c_i$.

By Definition \textcolor{blue}{\ref{defintion:NSR}} there exist $ p_0\geq 1$ and $0\leq j_0< k$ such that $I(a^0)=n_{p_0}+j_0$ and where
\begin{align}
    & a_{n_{p_0}}\ =\ c_1,\ \ldots,\ a_{n_{p_0} + j_0 -1}\ =\ c_{j_0},\  a_{n_{p_0}  + j_0}\ \geq \ c_{j_0+1}. \label{eq:3} 
\end{align}

Since $n_{p_0}+j_0$ is the first overfilled element of $a$, we have that $(a_n^0)_{n=1}^{n_{p_0} + j_0-1}$ is a $\vec{c}$-SR. We now proceed by case analysis to define $a^1$.\\ 

    \noindent \textbf{Case 1: $j_0=k-1$.} Then $(a_n^0)_{n=1}^{n_{p_0} + j_0}$ is end complete and so by applying Proposition \textcolor{blue}{\ref{prop:End complete}}, a finite sequence of carries transforms $(a_n^0)_{n=1}^{n_{p_0} + j_0}$ into a $\vtr{c}$-SR say $(b^{0}_n)_{n=1}^{n_{p_0} + j_0}$. Define $a^1$ by

    $$a_n^{1}\ =\ \begin{cases}
        b^0_n, & 1\leq n\leq n_{p_0}+j_0;\\
        a_n^0, & \text{otherwise.}
    \end{cases}$$

    \noindent In this case we have that $G(a^1)\leq G(a^0)-(c-1)$.\\

   \noindent \textbf{Case 2: $0\leq j_0<k-1$.} Note that in this case, $a_{n_{p_0}  + j_0}\ > \ c_{j_0+1}$. Borrow from $a_{n_{p_0}+j_0}^0$ to define $d^0$. In particular, define

    $$d^0_n\ =\ \begin{cases}
        a_n^0+c_i, & \text{if } n=n_{p_0}+j_0+i\ \text{for some }1\leq i\leq k;\\
        a_n^0-1, & \text{if } n=n_{p_0}+j_0;\\
        a_n^0 & \text{otherwise.}
    \end{cases}$$   
   
  \noindent Since $\vtr{c}$ is weakly decreasing, then for all $1\leq  i\leq k-j_0-1$ 
\begin{align*}
    &d^0_{n_{p_0} +j_0 +i}\ \geq\ c_{i}\ \geq\ c_{j_0+i+1}
 \end{align*}
and, thus,

\begin{align}
    &d^0_{n_{p_0} +j}\ \geq\ c_{j+1}\ \text{for each }0\leq j\leq k-1.
 \end{align}

\noindent Moreover, if we define $(e^0_n)_{n=1}^{n_{p_0} + k-1}$ by

$$e^0_n\ =\ \begin{cases}
        d^0_n, & \text{if  }1\leq n< n_{p_0}+j_0\ \text{ or }\ n_{p_0}+k\leq n;\\
        c_{i+1}, & \text{if  } n=n_{p_0}+i\ \text{ for some }j_0\leq i\leq k-1,
    \end{cases}$$

\noindent then $(e^0_n)_{n=1}^{n_{p_0} + k-1}$ is end complete and so by applying Proposition \textcolor{blue}{\ref{prop:End complete}}, a finite sequence of carries transforms $(e^0_n)_{n=1}^{n_{p_0} + k-1}$ into a $\vtr{c}$-SR, say $f^0:=(f^0_n)_{n=1}^{n_{p_0} + k-1}$. Define $a^1$ by

    $$a_n^{1}\ =\ \begin{cases}
        f^0_n, & 1\leq n< n_{p_0}+j_0;\\
        d^0_n-c_{i+1}, & \text{if } n=n_{p_0}+i\ \text{ for some }j_0\leq i\leq k-1;\\        d^0_n, & \text{otherwise.}
    \end{cases}$$

    \noindent In particular, $a^1$ is obtained from $a^0$ by a single borrow and then at least one carry. By Remark \textcolor{blue}{\ref{Rem:sumofcoefs bor/car}} this implies that $G(a^1)\leq G(d)-(c-1)=G(a^0)$.\\
    
    
If $a^1$ is defined in either case is a $\vtr{c}$-SR, then we are done. If not, then as above there exist $p_1\geq 1$ and $0\leq j_1< k$ such that $I(a^1)=n_{p_1}+j_1$ and where
\begin{align}
    & a_{n_{p_1}}\ =\ c_1,\ \ldots,\ a_{n_{p_1} + j_1 -1}\ =\ c_{j_1},\  a_{n_{p_1}  + j_1}\ >\ c_{j_1+1}. \label{eq:4} 
\end{align}

There are two important properties we prove.\\

\noindent \textbf{Claim 1:} We have that $n_{p_1}\geq n_{p_0}$ and if $G(a^1)=G(a^0)$ then $G_{n_{p_1}}(a^1)\geq G_{n_{p_0}}(a^0)+1$.\\

\begin{proof}[Proof of Claim 1.]
    If $a^1$ is constructed through Case 1 then $a^{1}_{n_{p_0}+k-1}=a^{0}_{n_{p_0}+k-1}-c_k=0$ and so the claim follows from the fact that $(b^0_n)_{n=1}^{n_{p_0}+j_0}=(a^1_n)_{n=1}^{n_{p_0}+j_0}$ is a $\vtr{c}$-SR. Note that in Case 1, we have $G(a^0)\neq G(a^1)$.\\

    Suppose instead that $a^1$ is constructed through Case 2. It is immediate that $n_{p_1}\geq n_{p_0}$ in the case where $n_{p_0}=1$ and so we may suppose further that $n_{p_0}\geq 2$. If $j_0=0$ then $n_{p_1}\geq n_{p_0}$ follows from the fact that $(f_n)_{n=1}^{n_{p_0}-1}=(a_n^1)_{n=1}^{n_{p_0}-1}$ is a $\vtr{c}$-SR. If $j_0\neq 0$ then $n_{p_1}\geq n_{p_0}$ follows from the fact that $a^{1}_{n_{p_0}+j_0-1}=a^{0}_{n_{p_0}+j_0-1}-c_{j_0}=0$ and since $(f_n)_{n=1}^{n_{p_0}+j_0}=(a_n^1)_{n=1}^{n_{p_0}+j_0}$ is a $\vtr{c}$-SR.

    Note that by Remark \textcolor{blue}{\ref{Rem:sumofcoefs bor/car}}, if $G(a^1)=G(a^2)$ then $a^1$ is obtained from $a^0$ by exactly one borrow from $a^0_{n_{p_0}}$ and then one carry into $a^{0}_{n_{p_0}-1}$ where $n_{p_0}-1\geq 1$. Thus, as single borrow from $a^0_{n_{p_0}}$ does not effect $G_{n_{p_0}}$ and a single carry into $a^{0}_{n_{p_0}-1}$ increases $G_{n_{p_0}}$ by 1, we have that

    $$G_{n_{p_1}}(a^1)\ \geq\  G_{n_{p_0}}(a^1)\ =\ G_{n_{p_0}}(a^0)+1$$

    \noindent as desired.
\end{proof}

Since $n_{p_1}+j_1$ is the first overfilled element of $a^1$, we have that $(a_n^1)_{n=1}^{n_{p_1} + j_1-1}$ is a $\vec{c}$-SR.

Inductively, if $a^{\ell}$ is a $\vtr{c}$-NSR and is defined for some $\ell\geq 1$, then we may define $n_{p_{\ell}}$ and $j_{\ell}$ such that $n_{p_\ell}+j_\ell$ is the first overfilled element of $a^\ell$. Then we have that $(a_n^\ell)_{n=1}^{n_{p_\ell} + j_\ell-1}$ is a $\vec{c}$-SR and employ the same case analysis as above to construct $a^{\ell+1}$. If $a^{\ell+1}$ is a $\vtr{c}$-SR then we are done. Otherwise, we may define $n_{p_{\ell+1}}$ and $j_{\ell+1}$ such that $n_{p_\ell+1}+j_{\ell+1}$ is the first overfilled element of $a^\ell$ and have that $(a_n^\ell)_{n=1}^{n_{p_\ell} + j_\ell-1}$ is a $\vec{c}$-SR. Furthermore, by applying Claim 1, $n_{p_{\ell+1}}\geq n_{p_{\ell}}$ and either $G(a^{\ell+1})<G(a^{\ell})$ or $G_{n_{p_{\ell+1}}}(a^{\ell+1})>G_{n_{p_{\ell}}}(a^{\ell})$.\\

\noindent \textbf{Claim 2:} There exists a $q\in \N$ such that $a^{q}$ is a $\vtr{c}$-SR.

\begin{proof}[Proof of Claim 2.]
    Suppose not. Let $\alpha=G(a^0)$, $\beta:=\lceil\frac{\alpha}{c-1}\rceil$, and $Z:=\{i\in \N\ :\ G(a^{i})<G(a^{i-1})\}$. Note that $|Z|\leq \beta$. Indeed, if $|Z|> \beta$ then there exists a subset $\{z_1,\ldots,z_{\beta+1}\}\subset Z$ where $z_{j}<z_{j+1}$ for each $1\leq j\leq \beta$. Then

    \begin{align*}
        G(a^{z_{\beta+1}})\ \leq\ G(a^{z_{\beta}})-(c-1)\ \leq\ G(a^{z_{\beta-1}})-2(c-1)\ &\leq\ \cdots\ \leq\ G(a^{z_1})-\beta(c-1)\\
        &\leq\ \alpha-(\beta+1)(c-1)\ <\ 0.
    \end{align*}

    Let $z:=\max\{i\in \N\ :\ i\in Z\}$ and define $\gamma:=G(a^{z})$. Note that for all $j>z$ it must be the case that $G(a^{j})=G(a^{j-1})$ and $G_{n_{p_j}}(a^j)\geq G_{n_{p_{j-1}}}(a^{j-1})+1$. However, if we examine $q=z+\gamma$ then

    $$\gamma\ <\ G_{n_{z}}(a^z)+\gamma\ \leq\ G_{n_{p_{q}}}(a^{q})\ \leq\  G(a^q)\ \leq\  G(a^z)\ =\ \gamma. $$

    Hence, we have the desired contradiction.
    
\end{proof}
The proof is complete as Claim 2 establishes that process must terminate in a finite number of iterations to a $\vtr{c}-SR$.\\
\end{proof}

\end{document}
